% META: {"id": "TSPE-COMBI-EX-021", "chapitre": "TSPE-COMBINATOIRE", "type_objet": "exercice", "capacites": ["TSPE-COMBINATOIRE-C4"], "capacites_codes": ["C4"], "methodes": ["M1"], "parcours": 1, "competences": ["raisonner", "communiquer"], "duree_min": 8, "mode_creation": "ex_nihilo", "sources_inspiration": [], "fichier_tex": "chapitres/TSPE-COMBINATOIRE/exercices/TSPE-COMBI-EX-021.tex", "corrige_tex": "chapitres/TSPE-COMBINATOIRE/corriges/TSPE-COMBI-CO-021.tex", "status": "approved"}
% BEGIN-VERIFY
% from sympy import binomial, Rational
% n, k = 8, 3
% lhs = binomial(n-1, k-1) + binomial(n-1, k)
% rhs = binomial(n, k)
% assert lhs == rhs
% END-VERIFY
\begin{exercice}{TSPE-COMBI-EX-007}{3}{18}
On considere un groupe $G$ de $8$ personnes, dont une personne particuliere nommee Alice. On souhaite denombrer les comites de $3$ personnes que l'on peut former a partir de $G$.
\begin{enumerate}
  \item Combien y a-t-il de comites de $3$ personnes contenant Alice ? (Justifier par un denombrement, sans formule factorielle.)
  \item Combien y a-t-il de comites de $3$ personnes ne contenant pas Alice ?
  \item En deduire, par le principe additif, le nombre total de comites de $3$ personnes parmi les $8$.
  \item Retrouver ce resultat directement, et verifier la coherence avec la relation de Pascal appliquee a $n=8$, $k=3$.
\end{enumerate}
\end{exercice}
